\documentclass[11pt,a4paper]{article}
\usepackage[utf8]{inputenc}
\usepackage[T1]{fontenc}
\usepackage{amsmath}
\usepackage{booktabs}
\usepackage{graphicx}
\usepackage{hyperref}
\usepackage[margin=2.5cm]{geometry}

\title{Market-basket analysis of actor co-occurrence\\in the IMDB Top 1000 movies}
\author{Parsa Hajian Nejad -- 48891A\\Master in Data Science for Economics, Universit\`a degli Studi di Milano\\
\small  a.y. 2025/26 -- Project 2}
\date{November 2026}

\begin{document}
\maketitle

\noindent Code: \url{https://github.com/Parsa-Hajian/amd-market-basket-imdb}

\section{Dataset and parts used}

The dataset is \textbf{IMDB Top 1000 Movies and TV Shows}, published on Kaggle by Harshit Shankhdhar.
It contains 1,000 rows and 15 columns including titles, ratings, genres, directors, and actors.
For this project, I have used only two pieces of information per row: the movie title (\texttt{Series\_Title}),
which identifies the basket, and the four credited actors (\texttt{Star1}--\texttt{Star4}), which are the items.
Other columns are not used for this project.

\section{Data organization}

Each movie is considered as a \textbf{basket} and its four credited actors as the \textbf{items} inside it.
The dataset generated 1,000 baskets: 996 of size~4 and 4 of size~3 (where one star column was empty or duplicated).
Across all baskets there are 2,709 distinct actors, but the distribution is extremely skewed:
2,068 actors (76\%) appear in a single movie, and only 271 actors appear in 3 or more movies.
The most frequent actors are Robert De Niro (17 movies), Tom Hanks (14), and Al Pacino (13).
A-Priori is a suitable approach for this amount of skew: pass~1 immediately eliminates most of the item space .

\section{Pre-processing}

Pre-processing is basic. For each row the four star columns are read; each name is stripped of leading
and trailing whitespace, and missing or empty values are dropped. The remaining names are collected into a
Python \texttt{set}, which removes any duplicates within the same movie, and then frozen into a
\texttt{frozenset}. A frozen set is immutable and hashable, which is required both for use as a dictionary
key (in the A-Priori counters) and for shipping baskets to Spark workers. A global flag \texttt{SUBSAMPLE}
allows the notebook to run on a smaller slice of the data during development.

\section{Algorithms and implementation}

\subsection{A-Priori (single machine)}

A-Priori exploits the \textbf{monotonicity} property: if an itemset is frequent (i.e.\ it appears in at
least $s$ baskets), then every subset of it is also frequent. Equivalently, if any subset is not frequent,
the itemset itself cannot be. The algorithm reads the baskets once per itemset size~$k$:

\begin{itemize}
  \item \textbf{Pass~1:} count every single actor. Keep those with count $\geq s$ (the ``frequent items'').
  \item \textbf{Pass~2:} for each basket, generate all pairs of its frequent items and increment a counter
        for each. Keep pairs with count $\geq s$.
  \item \textbf{Pass~$k$:} generate all $k$-subsets of a basket's frequent items; count only those whose
        every $(k{-}1)$-subset is already frequent (the prune test). Keep $k$-sets with count $\geq s$.
\end{itemize}

Counters are stored in a Python \texttt{Counter} (a hash table keyed by the itemset). This is the
``triples'' representation from the textbook: memory is proportional to the number of candidate itemsets
that actually occur in the data, not to the number of all possible pairs of frequent items. With baskets
of size~4, each basket contributes at most 6 pairs and 4 triples, so the number of occurring candidates
is small.

A brute-force implementation that counts every pair and triple without pruning is included in the notebook
as a correctness check: it must return the same frequent itemsets with the same counts.

\subsection{SON on Spark}

When baskets are too many to fit on one machine, the SON algorithm (Savasere, Omiecinski, Navathe)
splits them into chunks and processes each chunk independently. In this implementation, chunks correspond
to Spark partitions. SON runs in two rounds:

\begin{itemize}
  \item \textbf{Round~1 (\texttt{mapPartitions}):} each partition runs A-Priori locally with a scaled
        threshold $p \cdot s$, where $p$ is the fraction of baskets in that partition. Every locally
        frequent itemset is emitted as a candidate. The union of all partitions' candidates forms the
        global candidate set.
  \item \textbf{Round~2 (\texttt{flatMap} + \texttt{reduceByKey} + \texttt{filter}):} every basket is
        scanned once more; for each candidate it contains, a $(\text{candidate}, 1)$ pair is emitted.
        Counts are summed across all baskets, and only candidates with count $\geq s$ are kept.
\end{itemize}

The candidate set is broadcast to all workers so that round~2 does not require a shuffle for the lookup.

\textbf{No false negatives.} If an itemset has global count $\geq s$, it must have local count $\geq p \cdot s$
in at least one chunk (otherwise its global count would be below $s$). Therefore it will appear in the
candidate set. Round~2 removes the false positives.

\section{Scalability}

A-Priori on a single machine reads all baskets $k$ times (once per pass). Its memory is bounded by the
number of candidate itemsets that actually occur, which---with 4-item baskets---is small.

SON distributes the work: each executor needs only its own chunk in memory, plus the broadcast candidate
set (which is small). Both rounds are single-pass over the data, so the I/O cost scales linearly with the
number of baskets. The baskets are read from a text file with \texttt{sc.textFile}, which mirrors how data
would be loaded from distributed storage (HDFS, S3, GCS) in a production setting.

On a single machine with few cores, Spark's overhead (JVM startup, task scheduling, serialization)
dominates at small data sizes. The benefit appears when partitions are processed by many executors in
parallel on a cluster.

\section{Experiments}

\subsection{Effect of the support threshold}

Table~\ref{tab:threshold} shows how the number of frequent itemsets changes with the absolute support
threshold~$s$, on the full 1,000-basket dataset.

\begin{table}[h]\centering
\begin{tabular}{rrrr}\toprule
support $s$ & frequent actors & frequent pairs & frequent triples\\\midrule
2 & 641 & 121 & 25\\
3 & 271 & 25 & 3\\
4 & 145 & 5 & 1\\
5 & 79 & 3 & 1\\
6 & 59 & 1 & 0\\
8 & 17 & 0 & 0\\\bottomrule
\end{tabular}
\caption{Frequent itemsets as a function of the absolute support threshold (1,000 baskets).}
\label{tab:threshold}
\end{table}

At $s = 3$ the top frequent pairs are Radcliffe--Grint (6), Radcliffe--Watson (5), Watson--Grint (5),
Pesci--De Niro (4), and Allen--Hanks (4). The three frequent triples are the Harry Potter trio
(Radcliffe + Watson + Grint, support~5), the Star Wars trio (Fisher + Ford + Hamill, support~3),
and the Lord of the Rings trio (Wood + McKellen + Bloom, support~3).

\subsection{Scaling the number of baskets}

To study how the two implementations behave as the data grows, the baskets are replicated $m$ times and
shuffled, and the support threshold is scaled to $m \cdot s$ so that the relative support stays constant.
Table~\ref{tab:scaling} reports the results.

\begin{table}[h]\centering
\begin{tabular}{rrrrrr}\toprule
copies & baskets & $s$ & round-1 candidate pairs & A-Priori (s) & SON, 4 partitions (s)\\\midrule
1 & 1\,000 & 3 & 5\,833 & 0.00 & 1.55\\
10 & 10\,000 & 30 & 67 & 0.02 & 1.17\\
100 & 100\,000 & 300 & 25 & 0.18 & 1.74\\
500 & 500\,000 & 1\,500 & 25 & 0.82 & 4.79\\\bottomrule
\end{tabular}
\caption{Scaling the number of baskets at constant relative support (timings from a 1-core sandbox; update with Colab).}
\label{tab:scaling}
\end{table}

A-Priori time grows roughly linearly with the data size, as expected. SON time is dominated by Spark
overhead at small sizes and grows more slowly at large sizes, since each partition processes its chunk
independently. With 1 core the parallelism benefit is absent; on a multi-core or cluster setup, SON
would outperform A-Priori at larger scales. Note also that round-1 candidates drop sharply from 5,833
to 25 as the scaled threshold grows: with more data per chunk, fewer itemsets are false positives.

\subsection{Effect of the number of partitions}

Table~\ref{tab:partitions} shows how the number of chunks affects round-1 candidates and total SON time,
on 100,000 baskets with $s = 300$.

\begin{table}[h]\centering
\begin{tabular}{rrrrr}\toprule
partitions & baskets per chunk & $p\cdot s$ & round-1 candidates & SON (s)\\\midrule
1 & 100\,000 & 300.0 & 299 & 1.07\\
2 & 50\,000 & 150.0 & 299 & 1.28\\
4 & 25\,000 & 75.0 & 299 & 1.75\\
8 & 12\,500 & 37.5 & 324 & 2.65\\
16 & 6\,250 & 18.75 & 573 & 4.46\\\bottomrule
\end{tabular}
\caption{Effect of the number of chunks (100,000 baskets, $s=300$; 299 frequent itemsets in every case).}
\label{tab:partitions}
\end{table}

Smaller chunks mean lower local thresholds ($p \cdot s$), which produce more round-1 candidates (false
positives) that round~2 must recount. When $p \cdot s$ drops below~1, every itemset that occurs in a chunk
becomes a candidate and SON degenerates into a distributed brute-force count. On a single machine, more
partitions also mean more task-scheduling overhead, so time increases monotonically. On a cluster, the
optimal partition count balances parallelism against the false-positive cost.

\section{Discussion}

The project demonstrates the full A-Priori and SON pipeline on a small but real dataset.
The main findings are:

\begin{itemize}
  \item The IMDB Top 1000 dataset is dominated by actors who appear in a single movie, so A-Priori's
        pruning is very effective: pass~1 reduces 2,709 items to 271 at $s = 3$.
  \item Frequent pairs and triples are rare and correspond to well-known franchises (Harry Potter,
        Star Wars, Lord of the Rings) or long-running collaborations (De Niro--Pesci, Hanks--Allen).
  \item SON on Spark produces identical results to the single-machine baseline, as guaranteed by the
        no-false-negatives property.
  \item On a single core, Spark overhead makes SON slower than plain A-Priori. The benefit of SON
        would appear on a multi-core cluster where partitions are processed in parallel.
  \item The number of round-1 candidates is sensitive to the local threshold $p \cdot s$: more partitions
        means more false positives and a heavier round~2.
\end{itemize}

With only 1,000 baskets, the project demonstrates the machinery rather than producing statistically rich
patterns. On a real massive dataset (millions of transactions), the SON approach would be essential, and
the partition-count trade-off would become a meaningful tuning parameter.

\end{document}
